\documentclass[a4paper, exportdir=manual-output, exportenv=option,
    % uncomment the following line to enable export mode
    % exportimgs
]{tex2png}

\usepackage{enumitem}
\usepackage{microtype}
\usepackage{xurl}

\usepackage[hidelinks, pdfusetitle]{hyperref}

% adds item to description list for class option
% used instead of `exportimg`
\newenvironment{option}[2]{%
    \item[%
        \texttt{#1}%
        \ifx#2\empty\else=#2\fi%
    ]%
}{}

% disable footnotes in export mode
\ifexportimgs
    \makeatletter
    \let\footnote\@gobble
    \makeatother
\fi

\pdfstringdefDisableCommands{
    \let\,\relax
    \let\textsc\relax
}

\title{The \textsc{tex2png} class}
\author{G.\,P.\,S. Pennings}

\begin{document}

\maketitle

\begin{abstract}\noindent
    The \texttt{tex2png} class can be used to compile (parts of) a \LaTeX{} document to high-quality \textsc{png}%
    \footnote{Since \textsc{png} is widely supported, lossless, and supports transparency, it is arguably the best raster-graphics format for text and diagrams. Therefore, no other image formats are supported.}
    images, without the need to manually run external tools.
    These can be used, for example, in presentations or on the web when better alternatives are unavailable.

    A key feature of \texttt{tex2png} is the ability to quickly switch between two \emph{modes}:
    you can compile designated parts of the document on separate pages and automatically convert them to cropped images (\emph{export} mode),
    or you can compile the document normally to produce an overview of all content (\emph{overview} mode).

    The class is essentially a user-friendly wrapper around the \texttt{standalone} class, which, in turn, extends the standard \texttt{article}.
    This document is simultaneously the manual of \texttt{tex2png}, as well as an example of its use.
\end{abstract}



\section{Dependencies}
\label{sec:dependencies}

To convert each page of the \textsc{pdf} to an image, \texttt{tex2png} runs ImageMagick\footnote{\url{https://imagemagick.org/script/download.php}}\!, which in turns depends on Ghostscript\footnote{\url{https://ghostscript.com/releases/gsdnld.html}}\!.
These are both available on Windows, macOS (using Homebrew), and Linux.
This class was tested with ImageMagick~7.1.1--21 and Ghostscript~10.02.1, but other versions (of the same major version) will likely work as well.



\section{Quick start}
\label{sec:usage}

\begin{enumerate}
    \item Install the dependencies listed in \autoref{sec:dependencies};

    \item set your document's class to \texttt{tex2png} and enable the option \texttt{exportimgs};
    \begin{center}
        \verb@\documentclass[..., exportimgs]{tex2png}@
    \end{center}

    \item enclose the parts of your document that you want to convert to images in \texttt{exportimg} environments;
    % if no such environments are used, no PDF (or PNG) is generated
    \begin{center}
        \verb@\begin{exportimg} ... \end{exportimg}@
    \end{center}

    \item and compile your document with the \texttt{-shell-escape} flag\footnote{This is potentially dangerous, since it allows for arbitrary code execution by \LaTeX{} packages. However, this flag is required to use this class.}\!.
    \begin{center}
        \verb@pdflatex -shell-escape document.tex@
    \end{center}
\end{enumerate}

\noindent
Each page in the \textsc{pdf} has (excessively) wide margins, which are trimmed off by ImageMagick.
These margins ensure that, for instance, characters protruding from their boxes and list \verb@\item@'s are not cut off.%
\footnote{For an example, see \url{https://tex.stackexchange.com/q/53220}.}

Unfortunately, \texttt{tex2png} is not perfect: content may look different in export mode compared to overview mode.
Make sure to double check that the images look as intended.

% using this class to export figures is probably a bad idea: raster graphics will lose quality, and vector graphics should be kept as-is or converted directly



\section{Class options}
\label{sec:options}

This section lists all available options in \texttt{tex2png}.
Unrecognized options are passed to \texttt{article}.\footnote{Actually, if export mode is enabled, the options are passed to \texttt{standalone}, which then passes them to \texttt{article}. It is not recommended, however, to change the behavior of \texttt{standalone}.}

\begin{description}[font=\normalfont, style=nextline]
    \begin{option}{exportimgs}{}
        This option enables export mode: the base class is changed to \texttt{standalone} and each \texttt{exportimg} environment is converted to an image (all content outside these environments is ignored).
        Otherwise, overview mode is active, with \texttt{article} as the base class.
    \end{option}

    \begin{option}{exportdir}{\textlangle\textit{path}\textrangle}
        By default, the images are created in the current directory. This option can be used to change the output directory for images.
        If the specified directory does not exist, it is first created using a simple call to \texttt{mkdir}. For paths other than a single subdirectory, bear in mind the differences in path separator and implementation of \texttt{mkdir} between OSs.
    \end{option}

    \begin{option}{exportenv}{\textlangle\textit{string}\textrangle}
        By default, the ``exportimg'' environment is used to denote which parts of the document should be converted to images. With this option, you can change which environment is used.
        You can use existing (e.g., ``equation'') or custom environments for this. If the environment is not defined in the preamble, \texttt{tex2png} defines it for you (without functionality).
    \end{option}

    \begin{option}{exportsupersample}{\textlangle\textit{float}\textrangle}
        To improve image quality, the \textsc{pdf} is \emph{supersampled}, i.e., rendered at a higher resolution and then scaled down. The default supersampling rate is~$2$, and it is recommended to leave it unchanged.
    \end{option}

    \begin{option}{exportscale}{\textlangle\textit{float}\textrangle}
        Ignoring supersampling, the \textsc{pdf} is rendered at $72x$ pixels per inch, where~$x$ is the scale specified by this option. The default is $x=2$.
    \end{option}

    \begin{option}{exportbg}{{[\textit{color}]}}
        This option can be used to set the image background color, which is trans\-parent by default.
        Any color supported by ImageMagick\footnote{\url{https://imagemagick.org/script/color.php}} is allowed.
        If this option is specified without a color, white will be used.
        The same color is used for the border.
        Note that if your document already has a background (e.g., using the \texttt{pagecolor} package), that will take precedence.
    \end{option}

    \begin{option}{exportborder}{\textlangle\textit{string}\textrangle}
        After trimming the \textsc{pdf} margins, by default, a border of $5$px is added on all sides. The border width can be customized using this option.\footnote{See ImageMagick for the range of possible values: \url{https://imagemagick.org/script/command-line-options.php\#border}.}
    \end{option}
\end{description}



% advanced usage:
% • restricted shell
% • \ifexportimgs
% • how to get correct \ref's

% future work:
% • option to customize base class
% • improve cropping of PDF pages so they can be used as vector graphics

\end{document}
